% Experiment report. Fixed sections: see experiments/TEMPLATE/report.tex and experiments/README.md.
% Build: npm run experiment:report -- 12-perturbation-probe
\documentclass{../polyexp}
\begin{document}
% generated by scripts/experiment-report.ts from experiment.json and models/<run>/ - do not edit
\providecommand{\ExpHeader}{}\renewcommand{\ExpHeader}{%
{\sffamily\footnotesize\color{expgrey}EXPERIMENT 12 \textperiodcentered\ TUNING \hfill 2026-11-25}\par
{\LARGE\bfseries\sffamily Perturbation probe: tune all knobs with a significance test}\hfill\verdictbox{exporange}{RUNNING}\par\smallskip
{\itshape Can nudging every knob at random in every game, and moving only the knobs whose up and down games differ significantly, find gains that hill climbing cannot see?}\par\medskip
\begin{tabular}{@{}p{22mm}p{\dimexpr\linewidth-26mm}@{}}
\fact{Change}{A new trainer. Each sample nudges all knobs up or down by 25\% at random and plays the two opposite variations against each other and against a league; a knob moves half a nudge when its up-versus-down difference reaches z = 3.}
\fact{Knobs}{\knob{all non-zero knobs (about 200)}}
\fact{Base}{\model{greedy-full-combo}}
\fact{Method}{Tuning run of 100,000 games (10 rounds), against a hill-climbing control of the same budget; final test 400 games; metric: finish score}
\fact{Setups}{4 players, 16$\times$16, drylands, turn cap 150}
\fact{Record}{issue \issue{89}; status: running}
\end{tabular}\par\medskip
\begin{center}\begin{tikzpicture}[node distance=5mm]
    \node[flowbox] (n0) {{\sffamily\bfseries\scriptsize\color{expgrey}BASE}\\[1pt]\model{full-combo}};
    \node[flowbox,right=5mm of n0] (n1) {{\sffamily\bfseries\scriptsize\color{expgrey}CHANGE}\\[1pt]\knob{all non-zero knobs (about 200)}};
    \node[flowbox,right=5mm of n1] (n2) {{\sffamily\bfseries\scriptsize\color{expgrey}MEASURED}\\[1pt]100,000-game run vs a hill-climbing control\\{\scriptsize 1 setup, metric: finish score}};
    \node[flowbox,right=5mm of n2] (n3) {{\sffamily\bfseries\scriptsize\color{expgrey}RESULT}\\[1pt]--};
    \node[flowbox,right=5mm of n3] (n4) {{\sffamily\bfseries\scriptsize\color{expgrey}VERDICT}\\[1pt]\textbf{\color{exporange}RUNNING}};
    \draw[flowarrow] (n0) -- (n1); \draw[flowarrow] (n1) -- (n2); \draw[flowarrow] (n2) -- (n3); \draw[flowarrow] (n3) -- (n4);
\end{tikzpicture}\end{center}

}
\providecommand{\ExpResults}{}\renewcommand{\ExpResults}{%
\subsection{Perturbation probe, 100,000 games {\normalfont\small(\model{greedy-full-probe})}}
\begin{center}\begin{tikzpicture}\begin{axis}[expplot,height=4.6cm,xlabel={training games (thousands)},ylabel={Finish score, panel mean (\%)},enlarge y limits=0.2]
\addplot[dashed,expgrey,line width=0.7pt,forget plot] coordinates {(0.0,79.60) (0.0,79.60)};
\addplot[expblue,mark=*,mark size=1.8pt] coordinates {(0.0,79.60)};
\end{axis}\end{tikzpicture}\\[-2pt]{\footnotesize\color{expgrey}Learning curve of the tuning run (0 kept changes in 0 rounds). Dashed: the start.}\end{center}

\begin{center}\begin{tikzpicture}\begin{axis}[expplot,height=4.8cm,xmin=-5,xmax=5,ymin=0,xlabel={z-score of a knob (up games minus down games, in standard errors)},ylabel={knobs}]
\addplot[ybar interval,fill=expblue!55,draw=white] coordinates {(-5.0,3) (-4.5,0) (-4.0,0) (-3.5,2) (-3.0,4) (-2.5,5) (-2.0,13) (-1.5,19) (-1.0,24) (-0.5,37) (0.0,33) (0.5,28) (1.0,20) (1.5,8) (2.0,1) (2.5,1) (3.0,1) (3.5,0) (4.0,0) (4.5,6) (5.0,0)};
\addplot[dashed,expgrey,line width=0.8pt,domain=-5:5,samples=80] {205*0.5*exp(-x^2/2)/sqrt(2*pi)};
\draw[expgrey,line width=0.4pt] (axis cs:-3,0) -- (axis cs:-3,\pgfkeysvalueof{/pgfplots/ymax}) (axis cs:3,0) -- (axis cs:3,\pgfkeysvalueof{/pgfplots/ymax});
\end{axis}\end{tikzpicture}\\[-2pt]{\footnotesize\color{expgrey}z-scores of all 205 knobs after round 1 (10k training games). Dashed: pure chance. Vertical lines: the threshold $|z|=3$; 12 beyond it.}\end{center}
\begin{center}\sffamily\small\begin{tabular}{@{}rlrrr@{}}\toprule
Round & Knob & From & To & z\\\midrule
1 & \texttt{greedy.research} & 27.6 & 31.275 & $+7.6$\\
1 & \texttt{type.attack} & 1 & 1.133 & $+3.2$\\
1 & \texttt{tech.riding} & 0.948 & 0.837 & $-5.0$\\
1 & \texttt{tech.strategy} & 0.942 & 0.831 & $-3.1$\\
1 & \texttt{tech.climbing} & 1 & 1.133 & $+5.9$\\
1 & \texttt{tech.meditation} & 1.27 & 1.439 & $+6.6$\\
1 & \texttt{tech.philosophy} & 2.045 & 2.317 & $+5.8$\\
1 & \texttt{unit.warrior} & 1 & 1.133 & $+4.7$\\
1 & \texttt{unit.rider} & 1 & 0.882 & $-6.8$\\
1 & \texttt{reward.park} & 1 & 0.882 & $-7.3$\\
1 & \texttt{reward.superUnit} & 1 & 1.133 & $+7.6$\\
1 & \texttt{centre} & 1 & 0.882 & $-3.0$\\
\bottomrule\end{tabular}\\[3pt]{\footnotesize\color{expgrey}Knobs moved: 12.}\end{center}
\begin{center}\sffamily\small\begin{tabular}{@{}lrrrrrr@{}}\toprule
Knob & Value & Samples & Effect (pts) & z & z self & z league\\\midrule
\texttt{greedy.research} & 27.6 & 1667 & $+8.89$ & $+7.6$ & $+6.3$ & $+5.9$\\
\texttt{reward.superUnit} & 1 & 1667 & $+8.89$ & $+7.6$ & $+8.1$ & $+4.4$\\
\texttt{reward.park} & 1 & 1667 & $-8.59$ & $-7.3$ & $-8.2$ & $-3.9$\\
\texttt{unit.rider} & 1 & 1667 & $-7.96$ & $-6.8$ & $-6.3$ & $-4.7$\\
\texttt{tech.meditation} & 1.27 & 1667 & $+7.76$ & $+6.6$ & $+6.0$ & $+4.7$\\
\texttt{tech.climbing} & 1 & 1667 & $+6.89$ & $+5.9$ & $+5.5$ & $+4.0$\\
\texttt{tech.philosophy} & 2.045 & 1667 & $+6.81$ & $+5.8$ & $+4.5$ & $+4.8$\\
\texttt{tech.riding} & 0.948 & 1667 & $-5.84$ & $-5.0$ & $-4.4$ & $-3.6$\\
\texttt{unit.warrior} & 1 & 1667 & $+5.51$ & $+4.7$ & $+3.6$ & $+3.9$\\
\texttt{type.attack} & 1 & 1667 & $+3.82$ & $+3.2$ & $+3.4$ & $+1.9$\\
\bottomrule\end{tabular}\\[3pt]{\footnotesize\color{expgrey}Strongest evidence at the end of the run. Effect: finish score gained by one nudge up against one nudge down, halved.}\end{center}

\subsection{Control: hill climbing, 100,000 games {\normalfont\small(\model{greedy-full-probe-ctrl})}}
{\small\color{expgrey}No data recorded for this run yet.}
}

\ExpHeader

\section{Question and hypothesis}
Jakob, 25 November: ``Could we also add noise to other weights and then analyse the distribution and if we find a statistically significant difference between the positive and negative parts of the noise then we adjust the weights in that direction?''

Hill climbing accepts or rejects a whole bundle of changes on 300 to 400 games and found almost nothing in three runs of 100,000 games (experiments 1 to 3). The evolution-strategies trainer moves every weight on every step, significant or not. \textbf{Hypothesis:} measuring every knob in the same games and moving only the significant ones finds real single-knob gains that the other two trainers miss.

\textbf{Prediction, written before the run.} (1) Between 2 and 8 knobs move in 100,000 games, mostly research and tech weights (a 1,200-game smoke test already showed \knob{greedy.research} and three tech weights above $|z|=3$). (2) The panel finish score rises by 0.5 to 1.5 points. (3) Head to head against the base the result is 50 to 54\,\%, probably not clear of 50\,\% at $\pm$5. \textbf{Counts against the idea:} the z-scores of all knobs look like pure noise (spread near 1, nothing repeats after a knob's evidence is reset), or the panel curve falls as knobs move.

\section{Mechanism}
Each sample draws a sign $s_k\in\{-1,+1\}$ for every knob $k$, independently, and builds two opposite variations of the centre $\theta$: $w^{\pm}_k=\theta_k\,e^{\pm\sigma s_k}$ with $\sigma=0.25$ (a knob at 0 stays 0 and is not probed). On one map, every game from both seats, $w^+$ plays $w^-$, and both play the same league opponent. The outcome is
\[ x = \underbrace{\mathrm{finish}(w^+\text{ vs }w^-)-\tfrac12}_{\text{self}} \;+\; \underbrace{\mathrm{finish}(w^+\text{ vs }o)-\mathrm{finish}(w^-\text{ vs }o)}_{\text{league}} . \]
Because the signs are independent, the least-squares slope of $x$ on $s_k$ is $\hat\beta_k=\frac1n\sum_i x_i s_{ik}$: half the difference between the games where the knob was nudged up and those where it was nudged down. Its standard error is $\mathrm{sd}(x s_k)/\sqrt n$ and $z_k=\hat\beta_k/\mathrm{SE}$. After each round a knob with $|z_k|\ge 3$ moves half a nudge in the better direction and its evidence is reset; all other knobs keep accumulating evidence. Code: \texttt{scripts/train-probe.ts}. Course link: least squares and its analytic solution [EMAI ML-13, ML 2.0 p.27].

\section{Setup}
Base: the combined champion (economy plan plus movement noise 2), which beat both of its parents. 10 rounds of 1,667 samples (6 games each), league \model{hill-v4}, \model{full-v1} and the base. A curve point (200 games against each of the seven panel models, the same maps as every sweep) every two rounds; the centre of the best curve point is saved and plays the final test.

Sizing: one sample's outcome has a spread of about 48 points (smoke test), so after 100,000 games a knob that never moved has a standard error of about 0.4 points and is detected from about 1.1 points. With about 200 knobs and $|z|\ge 3$, about half a knob per round crosses the threshold by chance; such a knob has no real effect, so moving it should cost little, and the final test on unseen maps is the guard.

Control: hill climbing (\texttt{train-full}) from the same base with the skill's settings and the same 100,000 games.

\section{Results}
\ExpResults

\todo{claude}{Read the curve, the moves and the z-score distribution once both runs have finished.}

\section{Behaviour}
\todo{claude}{The sensitivity map: which knobs have a measurable effect at all, and which do not.}
\todo{jakob}{What you saw in the live games, if you watched any.}

\section{Interpretation}
\todo{claude}{Was the prediction right? Why or why not.}
\todo{jakob}{Your reading: agree, disagree, what is missing?}

\section{Threats to validity}
\todo{claude}{Two or three concrete threats.}

\section{Verdict and next step}
\todo{claude}{Verdict by the keep rules, once both runs are in.}

\section{Notes for the paper}
\todo{claude}{Draft the one-sentence claim.}
\todo{jakob}{Is this worth a place in the paper? Where?}

\end{document}
